% =====================================================================
\chapter{Series de potencias y series de Taylor}
\label{tema:6}
% =====================================================================

Una serie de potencias es un «polinomio de grado infinito». Dentro de su intervalo de convergencia se comporta como un polinomio: puede derivarse e integrarse término a término. Las funciones elementales $e^{x}$, $\sin x$, $\cos x$, $\ln(1+x)$ y $(1+x)^{a}$ son sumas de series de potencias, sus series de Taylor, y esa representación permite calcularlas con la precisión deseada, obtener límites y definir el número $e$. Este tema estudia el radio de convergencia, las operaciones con series y los desarrollos fundamentales, con demostración de que convergen a la función. Se siguen los capítulos~20 y~24 de Spivak~\cite{spivak}, los capítulos~7 y~11 de Apostol~\cite{apostol}, las §6.1 a~6.4 de~\cite{openstax2}, las §10.4 y~10.5 de Strang~\cite{strang} y las sesiones~97 a~100 del curso 18.01SC~\cite{mit1801}.

\section{Series de potencias y radio de convergencia}

\begin{definicion}[Serie de potencias]\label{def:serie-potencias}
Una \emph{serie de potencias} centrada en $c\in\R$ es una serie de la forma $\sum_{n=0}^{\infty}a_n(x-c)^{n}$, donde $(a_n)$ es una sucesión de números reales y $x$ es una variable real. Su \emph{dominio de convergencia} es el conjunto de los $x$ para los que converge.
\end{definicion}

Toda serie de potencias converge en $x=c$, donde vale $a_0$. El resultado siguiente muestra que el dominio de convergencia es siempre un intervalo centrado en $c$.

\begin{lema}[Lema de Abel]\label{lem:abel}
Si $\sum a_n(x_0-c)^{n}$ converge para algún $x_0\neq c$, entonces $\sum a_n(x-c)^{n}$ converge absolutamente para todo $x$ con $\abs{x-c}<\abs{x_0-c}$.
\end{lema}

\begin{proof}
Por la \cref{prop:condicion-necesaria}, $a_n(x_0-c)^{n}\to0$, y en particular existe $M$ con $\abs{a_n(x_0-c)^{n}}\le M$ para todo $n$. Sea $r=\frac{\abs{x-c}}{\abs{x_0-c}}<1$. Entonces $\abs{a_n(x-c)^{n}}=\abs{a_n(x_0-c)^{n}}\,r^{n}\le M\,r^{n}$, y la serie de valores absolutos converge por comparación con la serie geométrica $\sum Mr^{n}$.
\end{proof}

\begin{teorema}[Radio de convergencia]\label{teo:radio}
Para cada serie de potencias $\sum a_n(x-c)^{n}$ existe $R\in[0,\infty]$, su \emph{radio de convergencia}, tal que la serie converge absolutamente si $\abs{x-c}<R$ y diverge si $\abs{x-c}>R$. Si existe $\lim_{n\to\infty}\Bigl|\frac{a_n}{a_{n+1}}\Bigr|$ (finito o infinito), ese límite es $R$.
\end{teorema}

\begin{proof}
Sea $R$ el supremo (posiblemente $0$ o $\infty$) del conjunto de los valores $\abs{x_0-c}$ para los que la serie converge en $x_0$. Si $\abs{x-c}<R$, por la definición de supremo hay un $x_0$ de convergencia con $\abs{x-c}<\abs{x_0-c}$, y el \cref{lem:abel} da la convergencia absoluta en $x$. Si $\abs{x-c}>R$, la serie no converge en $x$, por la definición de $R$.

Para la fórmula, sea $\rho=\lim\abs{a_n/a_{n+1}}$ y sea $x\neq c$. El cociente de dos términos consecutivos de $\sum\abs{a_n(x-c)^{n}}$ es $\bigl|\frac{a_{n+1}}{a_n}\bigr|\abs{x-c}\to\frac{\abs{x-c}}{\rho}$. Por el \cref{teo:cociente}, la serie converge absolutamente si $\abs{x-c}<\rho$; si $\abs{x-c}>\rho$, el cociente tiende a un número mayor que~1, el término general no tiende a~$0$ y la serie diverge. Luego $R=\rho$.
\end{proof}

El intervalo $(c-R,c+R)$ es el \emph{intervalo de convergencia}. En los extremos $x=c\pm R$ el teorema no dice nada, y cada caso se estudia aparte con los criterios del \cref{tema:5}.

\begin{ejemplo}\label{ej:radios}
(a) $\sum_{n\ge0}x^{n}$: $R=1$, diverge en $x=\pm1$; dominio $(-1,1)$.
(b) $\sum_{n\ge1}\frac{x^{n}}{n}$: $\bigl|\frac{a_n}{a_{n+1}}\bigr|=\frac{n+1}{n}\to1$, luego $R=1$. En $x=1$ es la serie armónica, divergente; en $x=-1$ es la armónica alternada cambiada de signo, convergente. Dominio $[-1,1)$.
(c) $\sum_{n\ge0}\frac{x^{n}}{n!}$: $\bigl|\frac{a_n}{a_{n+1}}\bigr|=n+1\to\infty$, luego $R=\infty$.
(d) $\sum_{n\ge0}n!\,x^{n}$: $\bigl|\frac{a_n}{a_{n+1}}\bigr|=\frac{1}{n+1}\to0$, luego $R=0$ y solo converge en $x=0$.
\end{ejemplo}

\section{Derivación e integración término a término}

\begin{teorema}[Derivación e integración de series de potencias]\label{teo:termino-a-termino}
Sea $f(x)=\sum_{n=0}^{\infty}a_n(x-c)^{n}$ con radio de convergencia $R>0$. Las series
\[
\sum_{n=1}^{\infty}n\,a_n(x-c)^{n-1}\qquad\text{y}\qquad\sum_{n=0}^{\infty}\frac{a_n}{n+1}(x-c)^{n+1}
\]
tienen el mismo radio $R$; $f$ es derivable en $(c-R,c+R)$, con $f'$ igual a la primera serie, y $\int_c^{x}f(t)\dif t$ es igual a la segunda para todo $x$ de ese intervalo.
\end{teorema}

La demostración usa la convergencia uniforme de la serie en los intervalos cerrados contenidos en $(c-R,c+R)$ y puede consultarse en~\cite[cap.~24]{spivak} o \cite[cap.~11]{apostol}. Aplicando el teorema repetidamente, $f$ tiene derivadas de todos los órdenes en el intervalo de convergencia.

\begin{corolario}[Unicidad de los coeficientes]\label{cor:coeficientes}
Si $f(x)=\sum_{n=0}^{\infty}a_n(x-c)^{n}$ con $R>0$, entonces $a_n=\frac{f^{(n)}(c)}{n!}$ para todo $n$. En particular, dos series de potencias que coinciden en un entorno de $c$ tienen los mismos coeficientes.
\end{corolario}

\begin{proof}
Derivando $k$ veces término a término, $f^{(k)}(x)=\sum_{n\ge k}n(n-1)\cdots(n-k+1)\,a_n(x-c)^{n-k}$. En $x=c$ todos los términos con $n>k$ se anulan y queda $f^{(k)}(c)=k!\,a_k$.
\end{proof}

\begin{ejemplo}[Series obtenidas de la geométrica]\label{ej:desde-geometrica}
(a) Derivando $\frac{1}{1-x}=\sum_{n\ge0}x^{n}$ en $(-1,1)$: $\frac{1}{(1-x)^{2}}=\sum_{n\ge1}n\,x^{n-1}$.
(b) Sustituyendo $x$ por $-t$: $\frac{1}{1+t}=\sum_{n\ge0}(-1)^{n}t^{n}$ para $\abs t<1$. Integrando entre $0$ y $x$:
\[
\ln(1+x)=\sum_{n=1}^{\infty}\frac{(-1)^{n+1}}{n}x^{n}=x-\frac{x^{2}}{2}+\frac{x^{3}}{3}-\cdots,\qquad\abs x<1 .
\]
(c) Sustituyendo $x$ por $-t^{2}$ e integrando: $\arctan x=\sum_{n\ge0}\frac{(-1)^{n}}{2n+1}x^{2n+1}$ para $\abs x<1$.
\end{ejemplo}

\section{Series de Taylor}

\begin{definicion}[Serie de Taylor]\label{def:serie-taylor}
Sea $f$ con derivadas de todos los órdenes en un intervalo abierto que contiene a $c$. Su \emph{serie de Taylor} en $c$ es $\sum_{n=0}^{\infty}\frac{f^{(n)}(c)}{n!}(x-c)^{n}$. Para $c=0$ se llama también \emph{serie de Maclaurin}. La suma parcial de grado $n$ es el polinomio de Taylor $T_nf(x)$, y $R_nf(x)=f(x)-T_nf(x)$ es el \emph{resto}.
\end{definicion}

Por el \cref{cor:coeficientes}, si $f$ es suma de alguna serie de potencias centrada en $c$, esa serie es su serie de Taylor. Pero la serie de Taylor de $f$ puede converger y no tener por suma $f$ (\cref{eje:6-no-analitica}). La herramienta para decidirlo es el teorema de Taylor del módulo~0.5.

\begin{teorema}[Taylor con resto de Lagrange]\label{teo:taylor-lagrange}
Sea $f$ con $n+1$ derivadas en un intervalo abierto $I$ que contiene a $c$. Para cada $x\in I$ existe $\xi$ entre $c$ y $x$ tal que
\[
R_nf(x)=\frac{f^{(n+1)}(\xi)}{(n+1)!}(x-c)^{n+1}.
\]
\end{teorema}

\begin{proposicion}\label{prop:taylor-converge}
La serie de Taylor de $f$ en $c$ converge a $f(x)$ si y solo si $R_nf(x)\to0$ cuando $n\to\infty$.
\end{proposicion}

\begin{proof}
Las sumas parciales de la serie son $T_nf(x)=f(x)-R_nf(x)$, que tienden a $f(x)$ si y solo si $R_nf(x)\to0$.
\end{proof}

\section{Los desarrollos fundamentales}

\begin{teorema}[Exponencial]\label{teo:exp}
Para todo $x\in\R$, $\displaystyle e^{x}=\sum_{n=0}^{\infty}\frac{x^{n}}{n!}=1+x+\frac{x^{2}}{2}+\frac{x^{3}}{6}+\cdots$.
\end{teorema}

\begin{proof}
Todas las derivadas de $e^{x}$ son $e^{x}$ y valen~1 en $0$, así que la serie del enunciado es la de Taylor. Por el \cref{teo:taylor-lagrange}, $\abs{R_n(x)}=\frac{e^{\xi}\abs{x}^{n+1}}{(n+1)!}\le e^{\abs x}\frac{\abs{x}^{n+1}}{(n+1)!}$. La serie $\sum\frac{\abs{x}^{m}}{m!}$ converge (\cref{ej:cociente}), luego su término general tiende a~$0$ (\cref{prop:condicion-necesaria}), y $R_n(x)\to0$. Se concluye por la \cref{prop:taylor-converge}.
\end{proof}

Del mismo modo, como las derivadas de $\sin$ y $\cos$ son $\pm\sin$ y $\pm\cos$, acotadas por~1, sus restos cumplen $\abs{R_n(x)}\le\frac{\abs x^{n+1}}{(n+1)!}\to0$, y para todo $x\in\R$ (\cref{eje:6-seno-dem})
\[
\sin x=\sum_{n=0}^{\infty}\frac{(-1)^{n}x^{2n+1}}{(2n+1)!}=x-\frac{x^{3}}{6}+\cdots,\qquad
\cos x=\sum_{n=0}^{\infty}\frac{(-1)^{n}x^{2n}}{(2n)!}=1-\frac{x^{2}}{2}+\cdots.
\]

\begin{teorema}[Logaritmo]\label{teo:ln}
Para todo $x\in(-1,1]$, $\displaystyle\ln(1+x)=\sum_{n=1}^{\infty}\frac{(-1)^{n+1}}{n}x^{n}$. En particular, $\displaystyle\ln2=1-\frac12+\frac13-\frac14+\cdots$.
\end{teorema}

\begin{proof}
Para $t\neq-1$ y $n\in\N$, la suma geométrica finita da la identidad exacta
\[
\frac{1}{1+t}=\sum_{k=0}^{n-1}(-t)^{k}+\frac{(-t)^{n}}{1+t}.
\]
Integrando entre $0$ y $x\in(-1,1]$,
$\ln(1+x)=\sum_{k=1}^{n}\frac{(-1)^{k+1}x^{k}}{k}+E_n(x)$, con $E_n(x)=(-1)^{n}\int_0^{x}\frac{t^{n}}{1+t}\dif t$. Si $0\le x\le1$, como $1+t\ge1$, $\abs{E_n(x)}\le\int_0^{x}t^{n}\dif t=\frac{x^{n+1}}{n+1}\le\frac{1}{n+1}$. Si $-1<x<0$, para $t\in[x,0]$ se tiene $1+t\ge1+x>0$, y $\abs{E_n(x)}\le\frac{1}{1+x}\int_x^{0}\abs{t}^{n}\dif t=\frac{\abs{x}^{n+1}}{(n+1)(1+x)}$. En ambos casos $E_n(x)\to0$.
\end{proof}

La demostración cubre el extremo $x=1$, que el \cref{teo:termino-a-termino} no alcanza, y da además una cota explícita del error. La serie no sirve para $x>1$, porque allí su radio de convergencia, $R=1$, ya no la hace converger, aunque $\ln(1+x)$ esté definido.

\begin{teorema}[Serie binomial]\label{teo:binomial}
Sea $a\in\R$. Para todo $x\in(-1,1)$,
\[
(1+x)^{a}=\sum_{n=0}^{\infty}\binom{a}{n}x^{n},\qquad\binom{a}{n}=\frac{a(a-1)\cdots(a-n+1)}{n!},\quad\binom{a}{0}=1 .
\]
Si $a\in\N_0$, la serie es finita y se reduce al binomio de Newton.
\end{teorema}

Si $a\notin\N_0$, ningún coeficiente es nulo, $\bigl|\binom{a}{n}\big/\binom{a}{n+1}\bigr|=\frac{n+1}{\abs{a-n}}\to1$ y el radio es~1. La demostración de que la suma es $(1+x)^{a}$ se pide en el \cref{eje:6-binomial-dem}. Por ejemplo, con $a=\frac12$: $\sqrt{1+x}=1+\frac x2-\frac{x^{2}}{8}+\frac{x^{3}}{16}-\cdots$.

\begin{table}[h]
\centering
\small
\renewcommand{\arraystretch}{1.25}
\begin{tabular}{@{}>{$\displaystyle}l<{$}>{$\displaystyle}l<{$}>{$}l<{$}@{}}
\toprule
\text{Función} & \text{Serie de Maclaurin} & \text{Validez} \\
\midrule
\frac{1}{1-x} & \sum_{n\ge0}x^{n} & (-1,1) \\
e^{x} & \sum_{n\ge0}\frac{x^{n}}{n!} & \R \\
\sin x & \sum_{n\ge0}\frac{(-1)^{n}x^{2n+1}}{(2n+1)!} & \R \\
\cos x & \sum_{n\ge0}\frac{(-1)^{n}x^{2n}}{(2n)!} & \R \\
\ln(1+x) & \sum_{n\ge1}\frac{(-1)^{n+1}x^{n}}{n} & (-1,1] \\
\arctan x & \sum_{n\ge0}\frac{(-1)^{n}x^{2n+1}}{2n+1} & [-1,1] \\
(1+x)^{a} & \sum_{n\ge0}\binom{a}{n}x^{n} & (-1,1) \\
\bottomrule
\end{tabular}
\caption{Desarrollos fundamentales. La validez en $x=\pm1$ de la serie del arcotangente se prueba en el \cref{eje:6-arctan}.}
\label{tab:desarrollos}
\end{table}

\section{El número \texorpdfstring{$e$}{e}}

El número $e$ es la base de la exponencial, $e=e^{1}$. Por el \cref{teo:exp},
\[
e=\sum_{n=0}^{\infty}\frac{1}{n!}=1+1+\frac12+\frac16+\frac{1}{24}+\cdots\approx2{,}718281828 .
\]

\begin{proposicion}\label{prop:e-limite}
$\displaystyle\lim_{n\to\infty}\Bigl(1+\frac1n\Bigr)^{n}=e$.
\end{proposicion}

\begin{proof}
Se escribe $\bigl(1+\frac1n\bigr)^{n}=\exp\bigl(n\ln(1+\frac1n)\bigr)$. Con $h=\frac1n$, $n\ln(1+\frac1n)=\frac{\ln(1+h)-\ln1}{h}\to(\ln)'(1)=1$ cuando $h\to0$. Por la continuidad de la exponencial, el límite es $e^{1}=e$.
\end{proof}

\begin{teorema}[El número $e$ es irracional]\label{teo:e-irracional}
Sea $s_q=\sum_{k=0}^{q}\frac{1}{k!}$. Para todo $q\in\N$, $0<e-s_q<\frac{1}{q!\,q}$. En consecuencia, $e\notin\Q$.
\end{teorema}

\begin{proof}
El resto es $e-s_q=\sum_{k>q}\frac{1}{k!}>0$. Sacando factor común $\frac{1}{q!}$,
\[
e-s_q=\frac{1}{q!}\Bigl(\frac{1}{q+1}+\frac{1}{(q+1)(q+2)}+\cdots\Bigr)<\frac{1}{q!}\sum_{j=1}^{\infty}\frac{1}{(q+1)^{j}}=\frac{1}{q!}\cdot\frac{1/(q+1)}{1-1/(q+1)}=\frac{1}{q!\,q},
\]
donde la desigualdad es estricta porque $(q+1)(q+2)>(q+1)^{2}$. Supóngase ahora que $e=\frac pq$ con $p,q\in\N$. Entonces $x=q!\,(e-s_q)=q!\frac pq-\sum_{k=0}^{q}\frac{q!}{k!}$ es un entero, pues cada $\frac{q!}{k!}$ con $k\le q$ es entero y $q!\frac pq=(q-1)!\,p$. Pero $0<x<\frac1q\le1$, y no hay enteros estrictamente entre $0$ y $1$. Esta contradicción prueba que $e$ es irracional.
\end{proof}

La cota del teorema también es práctica: con $q=10$, $s_{10}$ aproxima $e$ con error menor que $\frac{1}{10!\cdot10}\approx2{,}8\cdot10^{-8}$.

\begin{observacion}[Errores frecuentes]\label{obs:6-errores}
(a) Dar el radio de convergencia como si fuera el dominio, sin estudiar los extremos. (b) Usar un desarrollo fuera de su intervalo de validez: $\ln3=\ln(1+2)$ no se obtiene sumando la serie en $x=2$, que diverge. (c) Suponer que la serie de Taylor de una función infinitamente derivable converge a la función (\cref{eje:6-no-analitica}). (d) En la serie binomial con $a$ negativo o fraccionario, olvidar los signos de los factores $(a-k)$.
\end{observacion}

\begin{observacion}[Conexión con el curso]
Los desarrollos de este tema son la base de muchas aproximaciones del curso. En el módulo~4 se usa el polinomio de Taylor de segundo orden, primeros términos de la serie de Taylor. En el módulo~6, la entropía de la distribución geométrica es $\sum_{k\ge1}k\,2^{-k}=2$ bits, suma que se obtiene del \cref{ej:desde-geometrica}(a) con $x=\frac12$ (\cref{eje:6-derivar-geom}). En el módulo~11, las bibliotecas numéricas calculan $\ln(1+x)$ y $e^{x}-1$ para $x$ pequeño con funciones específicas (\texttt{log1p}, \texttt{expm1}): restar $1$ tras redondear $1+x$ destruye las cifras significativas, mientras que los primeros términos de las series, $x-\frac{x^{2}}{2}$ y $x+\frac{x^{2}}{2}$, conservan la precisión. En el módulo~1, el límite de la softmax con temperatura cuando $T\to\infty$ se entiende con la aproximación de primer orden $e^{z_i/T}\approx1+\frac{z_i}{T}$: todos los numeradores tienden a~1 y la distribución, a la uniforme.
\end{observacion}

\begin{lecturas}
\item Spivak~\cite{spivak}, cap.~20 «Approximation by Polynomial Functions» y cap.~24 «Uniform Convergence and Power Series».
\item Apostol~\cite{apostol}, cap.~7 «Polynomial Approximations to Functions» y cap.~11 «Sequences and Series of Functions».
\item OpenStax~\cite{openstax2}, §6.1 «Power Series and Functions», §6.2 «Properties of Power Series», §6.3 «Taylor and Maclaurin Series» y §6.4 «Working with Taylor Series».
\item Strang~\cite{strang}, §10.4 «The Taylor Series for $e^x$, $\sin x$, and $\cos x$» y §10.5 «Power Series».
\item MIT 18.01SC~\cite{mit1801}, unidad~5, parte~B, sesiones~97 «Power Series», 98--99 «Taylor's Series» y~100 «Operations on Power Series».
\end{lecturas}

\section*{Ejercicios del tema 6}
\addcontentsline{toc}{section}{Ejercicios del tema 6}

\begin{ejercicio}{1}{eje:6-vf}
Justifique si son verdaderas o falsas: (a) si una serie de potencias centrada en $0$ converge en $x=3$, converge en $x=-2$; (b) si converge en $x=3$, converge en $x=-3$; (c) toda serie de potencias converge en al menos un punto; (d) la serie de Taylor de $\ln(1+x)$ en $0$ permite calcular $\ln5$.
\end{ejercicio}

\begin{ejercicio}{2}{eje:6-radio}
Halle el radio y el dominio de convergencia (estudiando los extremos): (a) $\sum_{n\ge1}\frac{x^{n}}{n^{2}}$; (b) $\sum_{n\ge1}n\,x^{n}$; (c) $\sum_{n\ge1}\frac{(x-2)^{n}}{3^{n}\,n}$; (d) $\sum_{n\ge0}\frac{(-1)^{n}x^{2n}}{4^{n}}$.
\end{ejercicio}

\begin{ejercicio}{2}{eje:6-e2x}
(a) Escriba la serie de $e^{2x}$, compruebe que sus coeficientes son $\frac{f^{(n)}(0)}{n!}$ y dé una serie cuya suma sea $e^{2}$. (b) Escriba la serie de $\sin2x$, compruebe con ella que la tercera derivada en $0$ vale $-8$ y dé una serie cuya suma sea $\sin2$.
\fuente{adaptado de Strang~\cite{strang}, §10.4, ejercicios 1 y~2.}
\end{ejercicio}

\begin{ejercicio}{2}{eje:6-taylor-basicas}
Calcule, a partir de las derivadas en $0$, las series de Maclaurin de: (a) $\frac{1}{1+x}$; (b) $\cosh x=\frac{e^{x}+e^{-x}}{2}$; (c) $\ln(1+x)$. Compruebe (b) sumando las series de $e^{x}$ y $e^{-x}$, y (c) con el \cref{teo:ln}.
\fuente{adaptado de Strang~\cite{strang}, §10.4, ejercicios 4, 6 y~8.}
\end{ejercicio}

\begin{ejercicio}{2}{eje:6-derivar-geom}
(a) Derive dos veces la serie geométrica y calcule $\sum_{n\ge1}\frac{n}{2^{n-1}}$ y $\sum_{n\ge2}\frac{n(n-1)}{2^{n-2}}$. (b) Deduzca $\sum_{k\ge1}k\,2^{-k}$ (la entropía en bits de la distribución geométrica de parámetro $\frac12$, módulo~6) y, para $0<p<1$, $\sum_{k\ge1}k(1-p)^{k-1}p$. (c) Sustituyendo $x$ por $1-x$ en la serie geométrica e integrando, obtenga la serie de $\ln x$ centrada en $1$.
\fuente{adaptado de Strang~\cite{strang}, §10.1, ejercicios 4 y~5.}
\end{ejercicio}

\begin{ejercicio}{2}{eje:6-productos}
Obtenga los primeros términos no nulos indicados, operando con series conocidas: (a) tres de $\frac{\sin x}{x}$; (b) cuatro de $e^{x}\sin x$; (c) cuatro de $e^{-x^{2}}$ y, de ahí, cuatro de $\int_0^{x}e^{-t^{2}}\dif t$.
\fuente{adaptado de Strang~\cite{strang}, §10.4, ejercicio~28.}
\end{ejercicio}

\begin{ejercicio}{2}{eje:6-limites}
Calcule con desarrollos de Taylor: (a) $\lim_{x\to0}\frac{e^{x}-1-x}{x^{2}}$; (b) $\lim_{x\to0}\frac{x-\sin x}{x^{3}}$; (c) $\lim_{x\to0}\frac{\ln(1+x)-x}{x^{2}}$; (d) $\lim_{x\to0}\frac{1-\cos x}{x^{2}}$.
\end{ejercicio}

\begin{ejercicio}{2}{eje:6-cos1}
Calcule $\cos1$ con cinco decimales correctos usando la serie de Taylor. ¿Cuántos términos hacen falta según la cota de error de las series alternadas?
\fuente{adaptado de Strang~\cite{strang}, §10.3, ejercicio~31.}
\end{ejercicio}

\begin{ejercicio}{2}{eje:6-aproximar}
(a) Escriba los cuatro primeros términos de $(1+x)^{1/2}$ y de $(1+x)^{-1/2}$, y aproxime $\sqrt{1{,}1}$. (b) Aproxime $\ln1{,}2$ con cuatro términos y acote el error con el \cref{teo:leibniz}. (c) ¿Cuántos términos de $\sum\frac{1}{k!}$ garantizan, por el \cref{teo:e-irracional}, un error menor que $10^{-6}$ en el cálculo de $e$?
\end{ejercicio}

\begin{ejercicio}{3}{eje:6-seno-dem}
Demuestre que $\sin x=\sum_{n\ge0}\frac{(-1)^{n}x^{2n+1}}{(2n+1)!}$ y $\cos x=\sum_{n\ge0}\frac{(-1)^{n}x^{2n}}{(2n)!}$ para todo $x\in\R$. Compruebe después, derivando término a término, que la derivada de la primera serie es la segunda.
\end{ejercicio}

\begin{ejercicio}{3}{eje:6-arctan}
(a) Integrando la identidad $\frac{1}{1+t^{2}}=\sum_{k=0}^{n-1}(-t^{2})^{k}+\frac{(-t^{2})^{n}}{1+t^{2}}$, demuestre que la serie del arcotangente converge a $\arctan x$ para todo $x\in[-1,1]$. (b) Deduzca la fórmula de Leibniz $\frac\pi4=1-\frac13+\frac15-\cdots$. (c) ¿Cuántos términos garantizan un error menor que $10^{-3}$ en $\frac\pi4$?
\end{ejercicio}

\begin{ejercicio}{3}{eje:6-binomial-dem}
Sea $a\in\R$ y $f(x)=\sum_{n\ge0}\binom{a}{n}x^{n}$ para $\abs{x}<1$. (a) Usando la identidad $(n+1)\binom{a}{n+1}+n\binom{a}{n}=a\binom{a}{n}$, demuestre que $(1+x)f'(x)=a\,f(x)$. (b) Demuestre que $g(x)=f(x)(1+x)^{-a}$ tiene derivada nula y concluya que $f(x)=(1+x)^{a}$.
\end{ejercicio}

\begin{ejercicio}{3}{eje:6-limite-exp}
Demuestre que $\lim_{n\to\infty}\bigl(1+\frac xn\bigr)^{n}=e^{x}$ para todo $x\in\R$, adaptando la demostración de la \cref{prop:e-limite}. Deduzca $\lim_{n\to\infty}\bigl(1-\frac1n\bigr)^{n}=\frac1e$ e interprete el resultado: si se eligen $n$ elementos al azar con reemplazamiento entre $n$, la probabilidad de que un elemento fijo no salga nunca tiende a $\frac1e\approx0{,}368$.
\end{ejercicio}

\begin{ejercicio}{3}{eje:6-no-analitica}
Sea $f(x)=e^{-1/x^{2}}$ si $x\neq0$ y $f(0)=0$. (a) Demuestre por inducción que, para $x\neq0$, $f^{(n)}(x)=P_n(1/x)\,e^{-1/x^{2}}$ con $P_n$ un polinomio. (b) Usando que $y^{k}e^{-y^{2}}\to0$ cuando $\abs y\to\infty$, demuestre que $f^{(n)}(0)=0$ para todo $n$. (c) Concluya que la serie de Taylor de $f$ en $0$ converge en todo $\R$ y que su suma solo coincide con $f$ en $x=0$.
\end{ejercicio}
